% GENERATED FILE. Edit canonical lessons, then run npm run generate:books.
% canonical-source-hash: fnv1a64:57be4feb04643b86
% canonical-lessons: MR-C473-khodkar, MR-C473-miskil, MR-C473-vinodi, MR-C473-gammat, MR-C473-thatta-karne, MR-R473-review-466-473

\chapter{Mischief and Fun}
\label{ch:mr-mischief-and-fun}
\hlchaptermodality{473}

\begin{chapteropening}
\textbf{By the end of this chapter:} \emph{I can describe mischief, fun and teasing.}

Complete the last lesson of chapter 473: I can describe mischief, fun and teasing.
\end{chapteropening}

\section[\texorpdfstring{\mr{खोडकर} (khoḍkar)}{खोडकर (khoḍkar)}]{\mr{खोडकर} (khoḍkar) — naughty}

\label{lesson:MR-C473-khodkar}



\begin{quote}
Before the new one: say the Marathi for hasty, then the Marathi for impatient.
\end{quote}

\subsection*{You'll want to know: \mr{खोडकर}}
\textbf{\mr{खोडकर}} — \emph{khoḍkar} — \textquotedblleft{}naughty\textquotedblright{}.

\begin{quote}
\mr{खोडकर} \mr{मुलगा} \mr{वर्गात} \mr{सारखा} \mr{बोलत} \mr{असतो}.
\end{quote}

\begin{quote}
\emph{khoḍkar mulgā vargāt sārkhā bolat asto.}
\end{quote}

\begin{quote}
The naughty boy keeps talking in class.
\end{quote}

Built on \mr{खोड}, a bad habit or prank (not the tree trunk): \mr{खोडी} \mr{काढणे}, to play a prank. Unchanging.

\begin{cousinweb}
\textbf{\mr{खोडकर}} (\emph{khoḍkar}) is a Marathi adjective, \mr{खोड}, \textquotedblleft{}mischief\textquotedblright{}, with the ending \emph{-kar}; the origin of \mr{खोड} is not settled.
\end{cousinweb}

\begin{grammarlens}[title={What you've built}]
One more word for describing what you lived and felt.
\end{grammarlens}

\subsection*{Your turn}
\begin{itemize}
\raggedright
  \item \emph{Say it:} \emph{khoḍkar}
  \item \emph{Say it:} \emph{khoḍkar}, once more
  \item \emph{Say it:} \emph{khoḍkar mulgā vargāt sārkhā bolat asto.} — The naughty boy keeps talking in class.
  \item \emph{Recall:} say the Marathi for hasty, then the Marathi for impatient, then say \emph{khoḍkar} again
\end{itemize}

\begin{tcolorbox}[breakable,colback=teal!4,colframe=teal!35!black,title={Before you move on}]
What does \mr{खोडकर} mean? (\textquotedblleft{}Naughty\textquotedblright{}.) Say it once more.
\end{tcolorbox}
\section[\texorpdfstring{\mr{मिश्किल} (miśkil)}{मिश्किल (miśkil)}]{\mr{मिश्किल} (miśkil) — impish; playful}

\label{lesson:MR-C473-miskil}



\begin{quote}
Before the new one: say the Marathi for impatient, then the Marathi for naughty.
\end{quote}

\subsection*{You'll want to know: \mr{मिश्किल}}
\textbf{\mr{मिश्किल}} — \emph{miśkil} — \textquotedblleft{}impish; playful\textquotedblright{}.

\begin{quote}
\mr{आजोबा} \mr{मिश्किल} \mr{हसले}.
\end{quote}

\begin{quote}
\emph{ājobā miśkil hasle.}
\end{quote}

\begin{quote}
Grandfather smiled impishly.
\end{quote}

Of a smile or a remark with a twinkle: \mr{मिश्किल} \mr{नजर}.

\begin{cousinweb}
\textbf{\mr{मिश्किल}} (\emph{miśkil}) is a Marathi adjective whose origin is not settled.
\end{cousinweb}

\begin{grammarlens}[title={What you've built}]
One more word for describing what you lived and felt.
\end{grammarlens}

\subsection*{Your turn}
\begin{itemize}
\raggedright
  \item \emph{Say it:} \emph{miśkil}
  \item \emph{Say it:} \emph{miśkil}, once more
  \item \emph{Say it:} \emph{ājobā miśkil hasle.} — Grandfather smiled impishly.
  \item \emph{Recall:} say the Marathi for impatient, then the Marathi for naughty, then say \emph{miśkil} again
\end{itemize}

\begin{tcolorbox}[breakable,colback=teal!4,colframe=teal!35!black,title={Before you move on}]
What does \mr{मिश्किल} mean? (\textquotedblleft{}Impish; playful\textquotedblright{}.) Say it once more.
\end{tcolorbox}
\section[\texorpdfstring{\mr{विनोदी} (vinodī)}{विनोदी (vinodī)}]{\mr{विनोदी} (vinodī) — funny; witty}

\label{lesson:MR-C473-vinodi}



\begin{quote}
Before the new one: say the Marathi for naughty, then the Marathi for impish.
\end{quote}

\subsection*{You'll want to know: \mr{विनोदी}}
\textbf{\mr{विनोदी}} — \emph{vinodī} — \textquotedblleft{}funny; witty\textquotedblright{}.

\begin{quote}
\mr{आमच्या} \mr{मावशी} \mr{खूप} \mr{विनोदी} \mr{आहेत}.
\end{quote}

\begin{quote}
\emph{āmcyā māvśī khūp vinodī āhet.}
\end{quote}

\begin{quote}
Our aunt is very funny.
\end{quote}

From \mr{विनोद}, a joke. Unchanging.

\begin{cousinweb}
\textbf{\mr{विनोदी}} (\emph{vinodī}) is a Marathi adjective built on \mr{विनोद}, \textquotedblleft{}a joke\textquotedblright{}, a tatsama from Sanskrit \emph{vinoda}, \textquotedblleft{}driving away; diversion\textquotedblright{} (Monier-Williams).
\end{cousinweb}

\begin{grammarlens}[title={What you've built}]
One more word for describing what you lived and felt.
\end{grammarlens}

\subsection*{Your turn}
\begin{itemize}
\raggedright
  \item \emph{Say it:} \emph{vinodī}
  \item \emph{Say it:} \emph{vinodī}, once more
  \item \emph{Say it:} \emph{āmcyā māvśī khūp vinodī āhet.} — Our aunt is very funny.
  \item \emph{Recall:} say the Marathi for naughty, then the Marathi for impish, then say \emph{vinodī} again
\end{itemize}

\begin{tcolorbox}[breakable,colback=teal!4,colframe=teal!35!black,title={Before you move on}]
What does \mr{विनोदी} mean? (\textquotedblleft{}Funny; witty\textquotedblright{}.) Say it once more.
\end{tcolorbox}
\section[\texorpdfstring{\mr{गंमत} (gammat)}{गंमत (gammat)}]{\mr{गंमत} (gammat) — amusement; a bit of fun}

\label{lesson:MR-C473-gammat}



\begin{quote}
Before the new one: say the Marathi for impish, then the Marathi for funny.
\end{quote}

\subsection*{You'll want to know: \mr{गंमत}}
\textbf{\mr{गंमत}} — \emph{gammat} — \textquotedblleft{}amusement; a bit of fun\textquotedblright{}.

\begin{quote}
\mr{पावसात} \mr{भिजायला} \mr{खूप} \mr{गंमत} \mr{आली}.
\end{quote}

\begin{quote}
\emph{pāvsāt bhijāylā khūp gammat ālī.}
\end{quote}

\begin{quote}
Getting soaked in the rain was great fun.
\end{quote}

Feminine.

\mr{गंमत} \mr{येणे}, to be fun, with the dative; \mr{गंमत} \mr{म्हणून}, as a joke.

\begin{cousinweb}
\textbf{\mr{गंमत}} (\emph{gammat}) is a Marathi noun whose origin is not settled.
\end{cousinweb}

\begin{grammarlens}[title={What you've built}]
One more word for describing what you lived and felt.
\end{grammarlens}

\subsection*{Your turn}
\begin{itemize}
\raggedright
  \item \emph{Say it:} \emph{gammat}
  \item \emph{Say it:} \emph{gammat}, once more
  \item \emph{Say it:} \emph{pāvsāt bhijāylā khūp gammat ālī.} — Getting soaked in the rain was great fun.
  \item \emph{Recall:} say the Marathi for impish, then the Marathi for funny, then say \emph{gammat} again
\end{itemize}

\begin{tcolorbox}[breakable,colback=teal!4,colframe=teal!35!black,title={Before you move on}]
What does \mr{गंमत} mean? (\textquotedblleft{}Amusement; a bit of fun\textquotedblright{}.) Say it once more.
\end{tcolorbox}
\section[\texorpdfstring{\mr{थट्टा} \mr{करणे} (thaṭṭā karṇe)}{थट्टा करणे (thaṭṭā karṇe)}]{\mr{थट्टा} \mr{करणे} (thaṭṭā karṇe) — to tease; to joke}

\label{lesson:MR-C473-thatta-karne}



\begin{quote}
Before the new one: say the Marathi for funny, then the Marathi for amusement.
\end{quote}

\subsection*{You'll want to know: \mr{थट्टा} \mr{करणे}}
\textbf{\mr{थट्टा} \mr{करणे}} — \emph{thaṭṭā karṇe} — \textquotedblleft{}to tease; to joke\textquotedblright{}.

\begin{quote}
\mr{मित्रांनी} \mr{माझ्या} \mr{टोपीची} \mr{थट्टा} \mr{केली}.
\end{quote}

\begin{quote}
\emph{mitrānnī mājhyā ṭopīcī thaṭṭā kelī.}
\end{quote}

\begin{quote}
My friends teased me about my cap.
\end{quote}

Past, with -\mr{ने} (he …): \mr{त्याने} \mr{थट्टा} \mr{केली}, \emph{tyāne thaṭṭā kelī}.

The target takes the genitive: \mr{भावाची} \mr{थट्टा} \mr{करणे}.

\begin{cousinweb}
\textbf{\mr{थट्टा} \mr{करणे}} (\emph{thaṭṭā karṇe}) is a Marathi phrase built on \mr{थट्टा}, \textquotedblleft{}a joke\textquotedblright{}, which CDIAL lists under the reconstructed \emph{*ṭhaṭṭha}, \textquotedblleft{}joke\textquotedblright{} (CDIAL 5492), with \mr{करणे}.
\end{cousinweb}

\begin{grammarlens}[title={What you've built}]
One more word for describing what you lived and felt.
\end{grammarlens}

\subsection*{Your turn}
\begin{itemize}
\raggedright
  \item \emph{Say it:} \emph{thaṭṭā karṇe}
  \item \emph{Say it:} \emph{thaṭṭā karṇe}, once more
  \item \emph{Say it:} \emph{mitrānnī mājhyā ṭopīcī thaṭṭā kelī.} — My friends teased me about my cap.
  \item \emph{Recall:} say the Marathi for funny, then the Marathi for amusement, then say \emph{thaṭṭā karṇe} again
\end{itemize}

\begin{tcolorbox}[breakable,colback=teal!4,colframe=teal!35!black,title={Before you move on}]
What does \mr{थट्टा} \mr{करणे} mean? (\textquotedblleft{}To tease; to joke\textquotedblright{}.) Say it once more.
\end{tcolorbox}
\section[(dialogue)]{Looking back: from Peace of Mind to Mischief and Fun}

\label{lesson:MR-R473-review-466-473}



\begin{quote}
Say the Marathi for satisfaction and well-being.
\end{quote}

\subsection*{Your turn}
Say these aloud:

\begin{itemize}
\raggedright
  \item satisfaction, well-being, relaxed, to rest a while, to brighten up
  \item torment, to writhe, to ache with pity, to stay alive, to wobble
  \item soft-hearted, mild, tender-hearted, sensitive, helpful to others
  \item harsh, merciless, wicked, fierce, deceitful
  \item sly, hypocritical, naive, innocent, to dupe
  \item sensible, shrewd, intelligent, understanding, common sense
  \item stubborn, obstinate, moody, hasty, impatient
  \item naughty, impish, funny, amusement, to tease
\end{itemize}

\begin{tcolorbox}[breakable,colback=teal!4,colframe=teal!35!black,title={Before you move on}]
Satisfaction? (\textbf{\mr{समाधान}}, \emph{samādhān}.) Well-being? (\textbf{\mr{स्वास्थ्य}}, \emph{svāsthya}.) Relaxed? (\textbf{\mr{निवांत}}, \emph{nivānt}.) To rest a while? (\textbf{\mr{विसावणे}}, \emph{visāvṇe}.) To brighten up? (\textbf{\mr{खुलणे}}, \emph{khulṇe}.) Torment? (\textbf{\mr{यातना}}, \emph{yātnā}.) To writhe? (\textbf{\mr{तडफडणे}}, \emph{taḍphaḍṇe}.) To ache with pity? (\textbf{\mr{कळवळणे}}, \emph{kaḷvaḷṇe}.) To stay alive? (\textbf{\mr{जगणे}}, \emph{jagṇe}.) To wobble? (\textbf{\mr{डगमगणे}}, \emph{ḍagmagṇe}.) Soft-hearted? (\textbf{\mr{मृदू}}, \emph{mṛdū}.) Mild? (\textbf{\mr{सौम्य}}, \emph{saumya}.) Tender-hearted? (\textbf{\mr{हळवा}}, \emph{haḷvā}.) Sensitive? (\textbf{\mr{संवेदनशील}}, \emph{saṁvedanśīl}.) Helpful to others? (\textbf{\mr{परोपकारी}}, \emph{paropkārī}.) Harsh? (\textbf{\mr{कठोर}}, \emph{kaṭhor}.) Merciless? (\textbf{\mr{निर्दयी}}, \emph{nirdayī}.) Wicked? (\textbf{\mr{दुष्ट}}, \emph{duṣṭa}.) Fierce? (\textbf{\mr{उग्र}}, \emph{ugra}.) Deceitful? (\textbf{\mr{कपटी}}, \emph{kapṭī}.) Sly? (\textbf{\mr{लबाड}}, \emph{labāḍ}.) Hypocritical? (\textbf{\mr{ढोंगी}}, \emph{ḍhoṅgī}.) Naive? (\textbf{\mr{भोळा}}, \emph{bhoḷā}.) Innocent? (\textbf{\mr{निरागस}}, \emph{nirāgas}.) To dupe? (\textbf{\mr{ठकवणे}}, \emph{ṭhakavṇe}.) Sensible? (\textbf{\mr{शहाणा}}, \emph{śahāṇā}.) Shrewd? (\textbf{\mr{चतुर}}, \emph{catur}.) Intelligent? (\textbf{\mr{बुद्धिमान}}, \emph{buddhimān}.) Understanding? (\textbf{\mr{समजूतदार}}, \emph{samjūtdār}.) Common sense? (\textbf{\mr{अक्कल}}, \emph{akkal}.) Stubborn? (\textbf{\mr{हट्टी}}, \emph{haṭṭī}.) Obstinate? (\textbf{\mr{आडमुठा}}, \emph{āḍmuṭhā}.) Moody? (\textbf{\mr{लहरी}}, \emph{lahrī}.) Hasty? (\textbf{\mr{उतावळा}}, \emph{utāvḷā}.) Impatient? (\textbf{\mr{अधीर}}, \emph{adhīr}.) Naughty? (\textbf{\mr{खोडकर}}, \emph{khoḍkar}.) Impish? (\textbf{\mr{मिश्किल}}, \emph{miśkil}.) Funny? (\textbf{\mr{विनोदी}}, \emph{vinodī}.) Amusement? (\textbf{\mr{गंमत}}, \emph{gammat}.) To tease? (\textbf{\mr{थट्टा} \mr{करणे}}, \emph{thaṭṭā karṇe}.)
\end{tcolorbox}
